\podnaslov{Večrazsežne zvezne porazdelitve}

\begin{definicija}
  Slučajni vektor $\sv{X}$ ima \pojem{zvezno porazdelitev}, če obstaja
  nenegativna funkcija $f_{\sv{X}}(\sv{x})$, da za $A \subseteq \R^n$ velja
  \[
	P(\sv{X} \in A) = \int_A f_{\sv{X}}(\sv{x}) d\sv{x}.
  \]
  Funkciji $f_{\sv{X}}$ pravimo \pojem{gostota}.
\end{definicija}

\vprasanje{Definiraj gostoto porazdelitve.}

\begin{izrek}
  Naj bo $f_{\sv{X}}(\sv{x})$ gostota vektorja $\sv{X}$ in $m < n$.
  Privzemimo, da je funkcija
  \[
	(x_1, \ldots, x_n) \mapsto \int_{\R^{n-m}} f_{\sv{X}}(x_1, \ldots, x_n)
	dx_{m+1} \ldots d_{x_n}
  \]
  Riemannovo integrabilna (lahko tudi v izlimitiranem smislu) po vseh Jordanovo
  izmerljivih množicah.
  Potem je to funkcija gostote vektorja $\sv{X}' = (x_1, \ldots, x_m)$.
\end{izrek}

\begin{izrek}
  Slučajna vektorja $\sv{X}, \sv{Y}$ sta neodvisna natanko tedaj, ko je
  \[
	f_{\sv{X}, \sv{Y}}(\sv{x}, \sv{y}) = f_{\sv{X}}(\sv{x}) f_{\sv{Y}}(\sv{y})
  \]
  skoraj povsod.
\end{izrek}

\begin{proof}
  V desno:
  Velja
  \begin{gather*}
	P(X \in A, Y \in B) = \int_{A \times B} f_{\sv{X}, \sv{Y}}(\sv{x}, \sv{y})
	d\sv{x} d\sv{y}, \\
	P(X \in A) P(Y \in B)
	= \int_A f_{\sv{X}}(\sv{x}) d\sv{x} \int_B f_{\sv{Y}}(\sv{y}) d\sv{y}
	= \int_{A \times B} f_{\sv{X}} f_{\sv{Y}}.
  \end{gather*}
  Ker sta vektorja neodvisna, sta ti količini enaki, torej sta integrirani
  funkciji enaki skoraj povsod.

  V levo:
  Velja
  \begin{align*}
	P(\sv{X} \in A, \sv{Y} \in B)
	&= \int_{A \times B} f_{\sv{X}}(\sv{x}) f_{\sv{Y}}(\sv{y}) d\sv{x} d\sv{y} \\
	&= \int_A f_{\sv{X}}(\sv{x}) d\sv{x} \int_B f_{\sv{Y}}(\sv{y}) d\sv{y} \\
	&= P(\sv{X} \in A) P(\sv{Y} \in B).
  \end{align*}
\end{proof}

\vprasanje{Karakteriziraj neodvisnost zvezno porazdeljenih slučajnih vektorjev
  in dokaži karakterizacijo.}

\begin{izrek}
  Naj bo $f_{\sv{X}, \sv{Y}}(\sv{x}, \sv{y}) = f(\sv{x}) g(\sv{y})$ za
  nenegativni $f, g$.
  Potem sta $\sv{X}$ in $\sv{Y}$ neodvisni.
\end{izrek}

Dokaz je praktično enak kot pri podobnem izreku za diskretne spremenljivke, le
da pišemo integrale namesto vsot.